\documentclass{article}
\usepackage{minted}
\title{Listing code examples}
\begin{document}
	\begin{listing}[!ht]
		\inputminted{octave}{BitXorMatrix.m}
		\caption{Example from external file}
		\label{listing:1}
	\end{listing}
	
	\begin{listing}[!ht]
		\begin{minted}{c}
			#include <stdio.h>
			int main() {
				printf("Hello, World!"); /*printf() outputs the quoted string*/
				return 0;
			}
		\end{minted}
		\caption{Hello World in C}
		\label{listing:2}
	\end{listing}
	
	\begin{listing}[!ht]
		\begin{minted}{lua}
			function fact (n)--defines a factorial function
			if n == 0 then
			return 1
			else
			return n * fact(n-1)
			end
			end
			
			print("enter a number:")
			a = io.read("*number") -- read a number
			print(fact(a))
		\end{minted}
		\caption{Example from the Lua manual}
		\label{listing:3}
	\end{listing}
	\noindent\texttt{minted} makes a nice job of typesetting listings \ref{listing:1}, \ref{listing:2} and \ref{listing:3}.
	\renewcommand\listoflistingscaption{List of source codes}
	\listoflistings
\end{document}